\subsection{Frobenius--Schur indicator and Larsen's alternative}

Recall (LIE, VIII, §7, n° 6, def. 2) that an irreducible representation of G in a finite-dimensional complex vector space E is said to be of orthogonal type (resp. of symplectic type) if there exists a nonzero symmetric bilinear form on E invariant under G (resp. if there exists a nonzero alternating bilinear form on E invariant under G). It is said to be of complex type if there exists no nonzero bilinear form on E invariant under G.

When an irreducible representation in E is of orthogonal type (resp. symplectic type), the space of symmetric (resp. alternating) G-invariant bilinear forms on E has dimension 1 and every nonzero G-invariant bilinear form on E is separating.

A representation of orthogonal type is sometimes said to be of real type, and a representation of symplectic type is sometimes said to be of quaternionic type (cf. LIE, IX, App. II).

Let \(\pi\) be an irreducible representation of G in a complex vector space E. The three possibilities above are distinguished by the value of the quantity

\[
\mathrm{FS} (\pi) = \int_ {\mathrm{G}} \chi_ {\pi} (g ^ {2}) d \mu (g),
\]

which is called the Frobenius--Schur indicator of π. Indeed we have

\[
\operatorname{FS} (\pi) = \left\{ \begin{array}{l l} 1 & \text{if }\pi\text{ is of orthogonal type,} \\ - 1 & \text{if }\pi\text{ is of symplectic type,} \\ 0 & \text{if }\pi\text{ is of complex type.} \end{array} \right.
\]

(LIE, IX, App. 2, p. 103, prop. 3 and p. 105, prop. 4).

When G is a Lie group, one can compute \(\mathrm{FS}(\pi)\) using prop. 1 of LIE, IX, p. 69.

DEFINITION 3. --- Let \(\varrho\) be a finite-dimensional unitary representation of G in a Hilbert space E. Let \(k\) be a positive integer. We call the dimension of the subspace of G-invariant elements in the representation \(\overline{\varrho}^{\otimes k} \otimes \varrho^{\otimes k}\) the absolute moment of order \(2k\) of \(\varrho\), and denote it by \(\mathrm{M}_{2k}(\varrho)\) .

THEOREM 3 (Larsen's alternative). --- Suppose that G is infinite. Let π be a faithful unitary representation of G in a Hilbert space E of finite dimension ≥ 2.

\begin{enumerate}
\def\labelenumi{\alph{enumi})}
\item
  Suppose that the derived group of G is infinite. We have \(M_{4}(\pi) \geqslant 2\) with equality if and only if G contains SU(E);
\item
  Suppose that \(\dim(E) \geqslant 3\) , that \(\pi\) is of orthogonal or symplectic type, and that \(\dim(E) \neq 4\) if \(\pi\) is of orthogonal type. Then \(M_4(\pi) \geqslant 3\) with equality if and only if the complexified Lie algebra of G is equal to the Lie algebra of a separating G-invariant bilinear form on E. *This is the case if and only if the neutral component \(G_0\) of G is a maximal compact subgroup of the group of automorphisms of such a bilinear form.*
\end{enumerate}

We shall use in the proof the following lemmas.

Lemma 5. --- Let \(\pi\) be a unitary representation of G in a Hilbert space E of finite dimension. Let \((\varrho_{i})_{i\in\mathrm{I}}\) be a family of nonzero unitary representations of G and \((n_{i})_{i\in\mathrm{I}}\) a family of integers \(\geqslant 1\) such that the representation \(\overline{\pi}\otimes\pi\) of G (resp. the representation \(\pi\otimes\pi\) ) is isomorphic to the direct sum \(\bigoplus_{i\in\mathrm{I}}\varrho_{i}^{n_{i}}\) . Then we have

\[
\mathrm{M} _ {4} (\pi) \geqslant \sum_ {i \in \mathrm{I}} n _ {i} ^ {2}
\]

with equality if and only if the representations \(\varrho_{i}\) are irreducible and pairwise non-isomorphic.

Denote by \(\chi_{i}\) the character of \(\varrho_{i}\) for \(i\in I\) . We have

\[
| \chi_ {\pi} | ^ {2} = \chi_ {\overline {{\pi}} \otimes \pi} = \sum_ {i \in \mathrm{I}} n _ {i} \chi_ {i}, \quad (\text { resp. } \chi_ {\pi} ^ {2} = \chi_ {\pi \otimes \pi} = \sum_ {i \in \mathrm{I}} n _ {i} \chi_ {i}).
\]

By the definition and cor. 3 of V, p. 466, we have

\[
\mathrm{M} _ {4} (\pi) = \int_ {\mathrm{G}} | \chi_ {\pi} | ^ {4} d \mu ,
\]

whence, in both cases, the formula

\[
\mathrm{M} _ {4} (\pi) = \sum_ {i, j} n _ {i} n _ {j} \int_ {\mathrm{G}} \chi_ {i} \overline {{\chi}} _ {j} d \mu = \sum_ {i, j} n _ {i} n _ {j} \dim \mathrm{Hom} _ {\mathrm{G}} (\varrho_ {i}, \varrho_ {j})
\]

(cor. 4, b) of V, p. 459). Let \(i \in I\) . Since the representation \(\varrho_{i}\) is not zero, the space \(\mathrm{Hom}_{\mathrm{G}}(\varrho_{i}, \varrho_{i})\) is not zero, and we deduce the lower bound

\[
\mathrm{M} _ {4} (\pi) \geqslant \sum_ {i \in \mathrm{I}} n _ {i} ^ {2}.
\]

Moreover, since \(n_{i} \geqslant 1\) , there is equality if and only if we have \(\dim \operatorname{Hom}_{\mathrm{G}}(\varrho_{i}, \varrho_{j}) = 0\) if \(i \neq j\) and \(\dim \operatorname{Hom}_{\mathrm{G}}(\varrho_{i}, \varrho_{i}) = 1\) for every i. The second condition holds if and only if the representations \(\varrho_{i}\) are irreducible (loc. cit.). The first is then satisfied if and only if the representations \(\varrho_{i}\) are pairwise non-isomorphic, by Schur's lemma (V, p. 387, cor. 2).

If E is a module over a commutative ring A, we will call the symmetric square (resp. exterior square) of E the A-module \(S^{2}(E)\) (resp. \(\wedge^{2}E\) ).

Lemma 6. --- Let \(q\) be a separating bilinear form on a finite-dimensional complex vector space E of dimension \(\geqslant 3\) .

\begin{enumerate}
\def\labelenumi{\alph{enumi})}
\item
  If \(q\) is symmetric, then the adjoint representation of the orthogonal Lie algebra \(\mathfrak{so}(q)\) is isomorphic to the exterior square \(\wedge^2\mathrm{E}\) of the natural representation \(\mathfrak{so}(q) \to \mathfrak{gl}(\mathrm{E})\) ;
\item
  If q is alternating, then the adjoint representation of the symplectic Lie algebra \(\mathfrak{sp}(q)\) is isomorphic to the symmetric square \(\mathsf{S}^2 (\mathrm{E})\) of the natural representation \(\mathfrak{sp}(q)\to \mathfrak{gl}(\mathrm{E})\) .
\end{enumerate}

Endow C with the identity involutive automorphism. The bilinear form q is then \(\varepsilon\) -Hermitian (A, IX, § 3, no. 1, def. 1) with \(\varepsilon = 1\) in case a) and \(\varepsilon = -1\) in case b). For every \(u \in \text{End}(\text{E})\) , we denote by \({}^{t}u\) the adjoint of u with respect to q. We therefore have \(q(x, u(y)) = q(^{t}u(x), y)\) for all \((x, y) \in \text{E} \times \text{E}\) . Let v denote the unique automorphism of E \(\otimes\) E such that \(v(x \otimes y) = y \otimes x\) for all \((x, y) \in \text{E}^{2}\) .

There exists a unique linear map \(w\) from \(\mathrm{E} \otimes \mathrm{E}\) into \(\operatorname{End}(\mathrm{E})\) such that \(w(a \otimes b)\) is the linear map defined by \(x \mapsto q(a, x)b\) for all \((a, b, x) \in \mathrm{E}^3\) . The map \(w\) is an isomorphism. For every \(s \in \mathrm{E} \otimes \mathrm{E}\) , we have

\[
{ } ^ { t } ( w ( s ) ) = \varepsilon w ( v ( s ) ) .\tag{3}
\]

Indeed, for all \((x,y)\) and \((a,b)\) in \(\mathrm{E} \times \mathrm{E}\) , it follows that

\[
q (x, w (a \otimes b) y) = q (a, y) q (x, b) = \varepsilon q (b, x) q (a, y) = q (\varepsilon w (b \otimes a) x, y).
\]

Finally, denote by \(x \mapsto x^*\) the isomorphism from E to \(\mathrm{E}'\) deduced from \(q\) , that is, such that \(\langle x^*, y \rangle = q(x, y)\) for all \((x, y) \in \mathrm{E}^2\) .

With these notations established, let H denote the orthogonal (resp. symplectic) group of q. It is a Lie subgroup of \(\mathbf{GL}(\mathrm{E})\) whose Lie algebra h is the subspace of End(E) consisting of the \(u \in \operatorname{End}(\operatorname{E})\) such that \(^{t}u = -u\) (LIE, III, p. 146, cor. 1). Since the form q is H-invariant, it follows that

\[
\langle (g a) ^ {*}, b \rangle = q (g a, b) = \langle a, g ^ {- 1} b \rangle
\]

for all \(g \in H\) and every \((a, b) \in E \times E\) .

Equip End(E) with the adjoint representation Ad of H. The linear map \(w\) is a morphism of representations of H: indeed, for every \(g \in \mathrm{H}\) and every \((a, b, x) \in \mathrm{E}^3\) , we have

\[
\begin{array}{c} w (g   (a \otimes b)) (x) = \langle (g a) ^ {*}, x \rangle   g b = \langle a, g ^ {- 1} x \rangle   g b \\ = g   (\langle a, g ^ {- 1} x \rangle   b) = \mathrm{Ad} (g) (w (a \otimes b)) x, \end{array}
\]

whence the conclusion by linearity.

Since v is also a morphism of representations of H, the same holds for the linear map \(\theta = w - \varepsilon w \circ v\) ; this is therefore a morphism of h-modules.

Denote by \(F \subset E \otimes E\) the subspace of elements \(s \in E \otimes E\) such that \(v(s) = -\varepsilon s\) . If \(\varepsilon = -1\) , the restriction to F of the canonical map from \(E \otimes E\) into the symmetric square of E is an isomorphism of C-vector spaces (A, III, p. 72); if \(\varepsilon = 1\) , the restriction to F of the canonical map from \(E \otimes E\) into the exterior square of E is an isomorphism of C-vector spaces (loc. cit., p. 82).

The image of F under \(\theta\) is contained in \(\mathfrak{h}\) : indeed, for every \(s\in\mathrm{F}\) the formula (3) implies

\[
{ } ^ { t } \theta ( s ) = { } ^ { t } w ( s ) - \varepsilon { } ^ { t } w ( v ( s ) ) = \varepsilon w ( v ( s ) ) - w ( s ) = - \theta ( s ) .
\]

By definition of F, the restriction of the map \(\theta\) to F coincides with that of \(2w\) and the linear map \(\theta\) is therefore injective on F. Since \(\dim(F) = \dim(h)\) . From (A, III, p. 75, th. 1 and p. 87, cor. 1 and LIE, VIII, p. 192 and p. 201), we conclude that \(\theta\) induces, by passage to subspaces, an isomorphism of \(h\)-modules from F onto \(h\), whence the lemma.

Lemma 7. --- Let \(H_{2}\) be a real Lie group and let \(H_{1}\) be a closed subgroup of \(H_{2}\) . The complexified Lie algebra of \(H_{1}\) is identified with a subrepresentation of the adjoint representation of \(H_{1}\) on the complexified Lie algebra of \(H_{2}\) .

Indeed, the restriction to \(H_{1}\) of the adjoint representation of \(H_{2}\) on its Lie algebra identifies with the adjoint representation of \(H_{1}\) .

Lemma 8. --- Let E be a complex Hilbert space of finite dimension n.

\begin{enumerate}
\def\labelenumi{\alph{enumi})}
\item
  The groups \(\mathbf{SU}(\mathrm{E})\) and \(\mathbf{U}(\mathrm{E})\) are connected;
\item
  The derived group of SU(E) is equal to SU(E);
\item
  The derived subgroup of U(E) is equal to SU(E).
\end{enumerate}

We may assume that \(n \geqslant 1\) and that \(E = C^{n}\) .

Let A be the subgroup of \(\mathbf{U}(\mathrm{E})\) formed by diagonal matrices; it is homeomorphic to \(U^{n}\) and therefore connected (TG, VI, p. 11, cor. 2 and I, p. 83, prop. 8). Its intersection with \(\mathbf{SU}(\mathrm{E})\) is homeomorphic to \(U^{n-1}\) , hence it is also connected.

By th. 1 of IV, p. 149, the group \(\mathbf{U}(\mathrm{E})\) (resp. \(\mathbf{SU}(\mathrm{E})\) ) is the union of the connected subsets \(g\mathrm{Ag}^{-1}\) for \(g \in \mathrm{G}\) (resp. of the connected subsets \(g(\mathrm{A} \cap \mathbf{SU}(\mathrm{E}))g^{-1}\) ); since the identity element belongs to each of these sets, the space \(\mathbf{U}(\mathrm{E})\) (resp. \(\mathbf{SU}(\mathrm{E})\) ) is connected (TG, I, p. 81, prop. 2).

Let us prove \(b\) ). The assertion is true when \(n = 1\) since \(\mathbf{SU}(\mathrm{E})\) is then reduced to the identity element. Suppose therefore that \(n \geqslant 2\) . The Lie algebra of \(\mathbf{SU}(\mathrm{E})\) is then simple (LIE, IX, p. 20, § 3, n°4), so the derived group of \(\mathbf{SU}(\mathrm{E})\) has finite index in \(\mathbf{SU}(\mathrm{E})\) . From this we deduce the result since \(\mathbf{SU}(\mathrm{E})\) is connected by \(a\) ).

The assertion \(c\) follows from \(b\) since the derived group of \(\mathbf{U}(\mathrm{E})\) is contained in \(\mathbf{SU}(\mathrm{E})\) .

We now prove Theorem 3. Let n denote the dimension of the Hilbert space E. Since the representation \(\pi\) is faithful, one may assume that G is a compact subgroup of U(E). The group G is closed in the real Lie group U(E), hence is a compact real Lie group (LIE, III, §8, no. 2, th. 2). By hypothesis, G is infinite, hence of dimension \(\geqslant 1\) . Let g be its Lie algebra; it is nonzero.

Let us prove \(a\) ). Suppose that the derived group D(G) is infinite. We have the G-invariant decomposition \(\operatorname{End}(\mathrm{E}) = \mathbf{C}1_{\mathrm{E}} \oplus \mathfrak{sl}(\mathrm{E})\) . The representations of G on \(\mathbf{C}1_{\mathrm{E}}\) and on \(\mathfrak{sl}(\mathrm{E})\) are not isomorphic, since \(\dim \mathfrak{sl}(\mathrm{E}) \geqslant 2\) . By Lemma 5, we therefore have \(\mathrm{M}_4(\pi) \geqslant 2\) with equality if and only if the representation of G in \(\mathfrak{sl}(\mathrm{E})\) is irreducible. Now, the derived group of G is contained in \(\mathbf{SL}(\mathrm{E})\) , and the complexification of the Lie algebra of D(G) therefore identifies with a subspace of \(\mathfrak{sl}(\mathrm{E})\) , which is a subrepresentation of the representation of G in \(\mathfrak{sl}(\mathrm{E})\) (Lemma 7). This subrepresentation is nonzero, since D(G) is infinite. We therefore have \(\mathrm{M}_4(\pi) = 2\) if and only if \((\mathcal{D}\mathfrak{g})_{(\mathbf{C})} = \mathfrak{sl}(\mathrm{E})\) . Since D(G) is contained in \(\mathbf{SU}(\mathrm{E})\) , and since \(\mathbf{SU}(\mathrm{E}) = \mathrm{D}(\mathbf{SU}(\mathrm{E}))\) (Lemma 8), this condition is equivalent to \(\mathbf{SU}(\mathrm{E}) \subset \mathbf{G}\) .

For the proof of assertion \(b\) ), we resume the notation of LIE, VIII, §13, nos. 2--4.

Suppose that \(\pi\) is of symplectic type and that \(\dim(\mathrm{E}) \geqslant 3\) . Let \(q\) be a nonzero G-invariant alternating bilinear form on E; it is separating (LIE, IX, p. 103, App. 2, prop. 3) and G is a compact subgroup of the symplectic group \(\mathbf{Sp}(q)\) . Denote by \(q^{*} \in \wedge^{2}\mathrm{E}\) the element to which the alternating form \(q\) identifies. The complexified Lie algebra \(\mathfrak{g}_{(\mathbf{C})}\) of G is contained in \(\mathfrak{sp}(q)\) . Now, since \(\dim(\mathrm{E}) \geqslant 3\) , the representation of \(\mathfrak{sp}(q)\) in \(\mathrm{E} \otimes \mathrm{E}\) admits the direct sum decomposition

\[
\mathrm{E} \otimes \mathrm{E} = \mathsf {S} ^ {2} (\mathrm{E}) \oplus \wedge^ {2} \mathrm{E} = \mathsf {S} ^ {2} (\mathrm{E}) \oplus \mathrm{E} _ {2} \oplus \mathbf {C} q ^ {*},
\]

where \(E_{2}\) is the second fundamental representation of \(\mathfrak{sp}(q)\) (LIE, VIII, p. 202, §13, no. 3, (IV)). The representations \(E_{2}\) and \(C q^{*}\) of \(\mathfrak{sp}(q)\) are irreducible (loc. cit.) and the representation \(S^{2}(E)\) is nonzero.

By Lemma 5, we therefore have \(\mathrm{M}_{4}(\pi) \geqslant 3\) with equality if and only if the representations of G in \(\mathsf{S}^{2}(\mathrm{E})\) , \(E_{2}\) and \(Cq^{*}\) are irreducible and pairwise non-isomorphic.

Suppose that \(M_{4}(\pi)=3\) . Since the representation \(S^{2}(E)\) identifies with the adjoint representation of \(\mathfrak{sp}(q)\) (Lemma 6), it contains as a subrepresentation the complexification of the adjoint representation of G (Lemma 7). We therefore have \(\mathfrak{g}_{(\mathbf{C})}=\mathfrak{sp}(q)\) .

Conversely, suppose that \(\mathfrak{g}_{(\mathbf{C})} = \mathfrak{sp}(q)\) . The adjoint representation of \(\mathfrak{sp}(q)\) and the representation \(\mathrm{E}_2\) of \(\mathfrak{sp}(q)\) are irreducible, of dimensions \(n(n + 1)/2\) and \(n(n - 1)/2 - 1\) respectively (LIE, VIII, p. 202, §13, no. 3, (IV)), which are different and \(\geqslant 2\) since \(n \geqslant 3\) , whence \(\mathrm{M}_4(\pi) = 3\) .

Finally, suppose that \(\pi\) is of real type and that \(\dim(\mathrm{E}) \geqslant 3\) and \(\dim(\mathrm{E}) \neq 4\) . Let \(q\) be a nonzero symmetric bilinear form on \(\mathrm{E}\), invariant under G ; it is separating (LIE, IX, p. 103, App. 2, prop. 3) and \(\mathrm{G}\) is a compact subgroup of the orthogonal group \(\mathbf{O}(q)\) . Denote by \(q^{*} \in \mathbb{S}^{2}(\mathrm{E})\) the element to which \(q\) identifies.

The complexified Lie algebra \(\mathfrak{g}_{(\mathbf{C})}\) of G is contained in \(\mathfrak{so}(q)\) . The representation of \(\mathfrak{so}(q)\) in E \(\otimes\) E admits the direct sum decomposition

\[
\mathrm{E} \otimes \mathrm{E} = \wedge^ {2} \mathrm{E} \oplus \mathsf {S} ^ {2} (\mathrm{E}) = \wedge^ {2} \mathrm{E} \oplus \mathsf {S} _ {0} ^ {2} (\mathrm{E}) \oplus \mathbf {C} q ^ {*},
\]

where \(\mathsf{S}_{0}^{2}(\mathbf{E})\) is the orthogonal complement of \(q^{*}\) in \(\mathsf{S}^{2}(\mathbf{E})\) . These representations are of dimension at least 2 since \(\dim(\mathbf{E})\geqslant3\) . By Lemma 5, we have \(\mathrm{M}_{4}(\pi)\geqslant3\) with equality if and only if the representations of G in \(\wedge^{2}\mathrm{E}\) and \(\mathsf{S}_{0}^{2}(\mathbf{E})\) are irreducible and non-isomorphic.

Suppose that \(\mathrm{M}_{4}(\pi)=3\) . Since the representation \(\wedge^{2}E\) identifies with the adjoint representation of \(\mathfrak{so}(q)\) (Lemma 6), it contains as a subrepresentation the complexification of the adjoint representation of G (Lemma 7). The condition \(\mathrm{M}_{4}(\pi)=3\) therefore implies that \(\mathfrak{g}_{(\mathbf{C})}=\mathfrak{so}(q)\) .

Conversely, suppose that \(\mathfrak{g}_{(\mathbf{C})} = \mathfrak{so}(q)\) . The adjoint representation of \(\mathfrak{so}(q)\) is irreducible, since \(\dim(\mathrm{E}) \neq 4\) (LIE, VIII, § 13, p. 193, (I) and p. 206, (I)), and of dimension \(n(n-1)/2\) . The representation \(\mathsf{S}_0^2(\mathrm{E})\) of \(\mathfrak{so}(q)\) has as highest weight \(2\varpi_1\) . By comparing its dimension with that of the irreducible representation of \(\mathfrak{so}(q)\) of highest weight \(2\varpi_1\) , computed by the Weyl formula (cf. LIE, VIII, §9, no. 2, th. 2 and LIE, VI, plates II and IV), one verifies that \(\mathsf{S}_0^2(\mathrm{E})\) is irreducible. We therefore conclude that \(\mathrm{M}_4(\pi) = 3\) if \(\mathfrak{g}_{(\mathbf{C})} = \mathfrak{so}(q)\) .

Remarks. --- 1) The condition “G is infinite” is necessary in Theorem 3 (exercise 9 of V, p. 507).

\begin{enumerate}
\def\labelenumi{\arabic{enumi})}
\setcounter{enumi}{1}
\item
  Let \(n \in \mathbf{N}\) . For every unitary representation \(\varrho\) of a compact group in a Hilbert space of dimension \(n\) , we have \(\mathrm{M}_{4}(\varrho) \leqslant n^{2}\) ; equality is possible (exercise 15 of V, p. 508).
\item
  One can prove (cf. R. GURALNICK and P.H. TIEP, Decomposition of small tensor powers and Larsen's conjecture, Representation Theory 9 (2005), 138--208) that the condition that G is infinite can be omitted if one assumes that E has dimension \(\geqslant 7\) and if one replaces the hypothesis \(M_{4}(\pi) = 2\) (resp. \(M_{4}(\pi) = 3\) ) by \(M_{8}(\pi) = 24\) (resp. \(M_{8}(\pi) = 105\) ).
\end{enumerate}

